\uuid{M0kv}
\titre{Matrice d'une application linéaire}
\auteur{hanani ; mzari}
\organisation{exo7}
\licence{CC BY-NC-SA 4.0}
\datecreate{2020-01-20}
\niveau{L1}
\module{Algèbre}
\chapitre{Matrice}
\sousChapitre{Matrice et application linéaire}
\difficulte{1}
\theme{Matrice d'une application linéaire}
\exo7qcmid{lille-2.matrice-application.004}

\contenu{
\begin{question}
\qclassification{108.03}
\qtags{motcle=Matrice d'une application linéaire}
On considère $\Rr^2$ muni de sa base canonique et $f$ l'application linéaire définie par :
$$\begin{array}{rccc}f:&\Rr^2 &\to& \Rr^2\\
& (x,y)&\to &(2x+y,4x-3y).\end{array}$$
Quelles sont les assertions vraies ?
\begin{answers}
\bad{La matrice de $f$ dans la base canonique est : $
\left(\begin{array}{rc}
2&4\\
1&-3\\
\end{array}\right)$}.
\good{La matrice de $f$ dans la base canonique est : $
\left(\begin{array}{rc}
2&1\\
4&-3\\
\end{array}\right)$}.
\good{$f$ est injective}.
\good{ $f$ est bijective}.
\end{answers}
\begin{explanations} Soit ${\cal {B}}=\{e_1,e_2\}$ la base canonique de $\Rr^2$. La matrice de $f$ dans la base  ${\cal {B}}$ est la matrice  dont la $j$ième colonne est constituée des coordonnées de  $f(e_j)$ dans la base ${\cal {B}}$. Cette matrice est :
 $
\left(\begin{array}{rc}
2&1\\
4&-3\\
\end{array}\right)$.
\vskip0mm
On vérifie que $\ker f=\{(0,0)\}$, donc $f$ est un endomorphisme  injectif de $\Rr^2$,
et donc $f$ est bijectif.
\end{explanations}
\end{question}
}
